\documentclass[11pt]{article}

\usepackage[utf8]{inputenc}
\usepackage[spanish]{babel}
\usepackage{amsmath}
\usepackage{amssymb}
\usepackage{geometry}
\geometry{margin=2.5cm}

\title{Desarrollo de Series de Fourier: \\ Onda Cuadrada y Onda Diente de Sierra}
\author{Curso Básico de MATLAB}
\date{}

\begin{document}

\maketitle

\section{Introducción}

La serie de Fourier de una función periódica $f(t)$, con periodo
$T = 2L$, se define como

\begin{equation}
f(t) = \frac{a_0}{2} + \sum_{n=1}^{\infty} \left[ a_n \cos\!\left(\frac{n\pi t}{L}\right) + b_n \sin\!\left(\frac{n\pi t}{L}\right) \right]
\end{equation}

donde los coeficientes se calculan como

\begin{align}
a_0 &= \frac{1}{L} \int_{-L}^{L} f(t)\, dt \\
a_n &= \frac{1}{L} \int_{-L}^{L} f(t) \cos\!\left(\frac{n\pi t}{L}\right) dt \\
b_n &= \frac{1}{L} \int_{-L}^{L} f(t) \sin\!\left(\frac{n\pi t}{L}\right) dt
\end{align}

En ambos casos de este documento se trabaja con periodo $T = 2\pi$
(es decir, $L = \pi$), y ambas funciones son \textbf{impares}, por
lo que $a_0 = 0$ y $a_n = 0$ para toda $n$: solo es necesario
calcular los coeficientes $b_n$.

\section{Caso 1: Onda cuadrada}

\subsection{Definición de la función}

Se define la onda cuadrada de periodo $2\pi$ como

\begin{equation}
f(t) =
\begin{cases}
\phantom{-}1, & 0 < t < \pi \\
-1, & -\pi < t < 0
\end{cases}
\end{equation}

Esta función es impar y equivale a $\operatorname{sign}(\sin t)$.

\subsection{Cálculo de $b_n$}

Como $f(t)$ es impar, únicamente aparece el coeficiente $b_n$:

\begin{equation}
b_n = \frac{1}{\pi} \int_{-\pi}^{\pi} f(t) \sin(nt)\, dt
    = \frac{2}{\pi} \int_{0}^{\pi} (1)\sin(nt)\, dt
\end{equation}

(se usó que el integrando $f(t)\sin(nt)$ es una función par, por
lo que la integral en $[-\pi,\pi]$ es el doble de la integral en
$[0,\pi]$).

Resolviendo la integral:

\begin{align}
b_n &= \frac{2}{\pi} \left[ -\frac{\cos(nt)}{n} \right]_0^{\pi}
     = \frac{2}{n\pi} \Big[ 1 - \cos(n\pi) \Big]
\end{align}

Como $\cos(n\pi) = (-1)^n$, se tiene:

\begin{equation}
b_n = \frac{2}{n\pi}\left[1 - (-1)^n\right] =
\begin{cases}
\dfrac{4}{n\pi}, & n \text{ impar} \\[2mm]
0, & n \text{ par}
\end{cases}
\end{equation}

\subsection{Ecuación final a programar}

Sustituyendo $b_n$ en la serie y dejando solo los términos impares
($n = 1, 3, 5, \ldots$):

\begin{equation}
\boxed{\;f(t) = \frac{4}{\pi} \sum_{n=1,3,5,\ldots}^{\infty} \frac{\sin(nt)}{n}\;}
\end{equation}

Esta es la expresión implementada en \texttt{onda\_cuadrada\_fourier.m},
sumando un número finito $K$ de armónicos impares en lugar de
infinitos.

\section{Caso 2: Onda diente de sierra}

\subsection{Definición de la función}

Se define la onda diente de sierra de periodo $2\pi$ como

\begin{equation}
f(t) = t, \qquad -\pi < t < \pi
\end{equation}

repitiéndose con periodo $2\pi$. También es una función impar.

\subsection{Cálculo de $b_n$}

\begin{equation}
b_n = \frac{1}{\pi} \int_{-\pi}^{\pi} t \sin(nt)\, dt
    = \frac{2}{\pi} \int_{0}^{\pi} t \sin(nt)\, dt
\end{equation}

(de nuevo, $t\sin(nt)$ es par, así que se duplica la integral en
$[0,\pi]$). Se resuelve por partes, con $u = t$, $dv = \sin(nt)\,dt$,
por lo que $du = dt$ y $v = -\cos(nt)/n$:

\begin{align}
\int_0^{\pi} t \sin(nt)\, dt
&= \left[ -\frac{t\cos(nt)}{n} \right]_0^{\pi} + \frac{1}{n}\int_0^{\pi} \cos(nt)\, dt \\
&= -\frac{\pi \cos(n\pi)}{n} + \frac{1}{n}\left[\frac{\sin(nt)}{n}\right]_0^{\pi} \\
&= -\frac{\pi (-1)^n}{n} + 0
\end{align}

(el segundo término se anula porque $\sin(n\pi) = 0$ para toda $n$
entera). Sustituyendo de vuelta:

\begin{equation}
b_n = \frac{2}{\pi}\left(-\frac{\pi(-1)^n}{n}\right) = \frac{2(-1)^{n+1}}{n}
\end{equation}

\subsection{Ecuación final a programar}

\begin{equation}
\boxed{\;f(t) = 2 \sum_{n=1}^{\infty} \frac{(-1)^{n+1}}{n} \sin(nt)\;}
\end{equation}

A diferencia de la onda cuadrada, aquí participan \textbf{todos}
los armónicos ($n = 1, 2, 3, \ldots$), alternando de signo. Esta es
la expresión implementada en
\texttt{onda\_diente\_sierra\_fourier.m}, sumando un número finito
$N$ de términos en lugar de infinitos.

\section{Fenómeno de Gibbs}

Ambas funciones estudiadas tienen discontinuidades de salto (la onda
cuadrada salta entre $-1$ y $1$; la diente de sierra salta entre
$\pi$ y $-\pi$ en cada periodo). Al aproximar una función
discontinua con un número finito de términos de su serie de
Fourier, aparece una sobre-oscilación (\emph{overshoot}) justo
antes y después de cada salto que no desaparece al aumentar el
número de términos $N$: lo que ocurre es que dicha oscilación se
concentra cada vez en un intervalo más angosto alrededor de la
discontinuidad, pero su \textbf{altura converge a un valor fijo},
aproximadamente igual al $9\%$ del tamaño del salto. Este
comportamiento se conoce como el \textbf{fenómeno de Gibbs} y puede
observarse directamente en las gráficas generadas por
\texttt{onda\_cuadrada\_fourier.m} y
\texttt{onda\_diente\_sierra\_fourier.m}: sin importar cuántos
armónicos se agreguen, siempre queda un pequeño pico cerca de cada
discontinuidad.

Formalmente, si $S_N(t)$ es la suma parcial con $N$ términos de la
serie de Fourier y $t_0$ es un punto de discontinuidad con salto de
tamaño $d = f(t_0^+) - f(t_0^-)$, se cumple que

\begin{equation}
\lim_{N \to \infty} \left( \max_{t \,\text{cerca de}\, t_0} S_N(t) - f(t_0^+) \right) \approx 0.0895 \, d
\end{equation}

es decir, el sobrepaso máximo no tiende a cero, sino a
aproximadamente el $9\%$ del salto, independientemente de cuántos
términos se utilicen.

\section{Fenómeno de Aliasing}

El aliasing es un fenómeno distinto al de Gibbs: no ocurre al
truncar una serie de Fourier, sino al \textbf{muestrear} (discretizar)
una señal continua para representarla con un número finito de
puntos, como sucede en cualquier sistema digital.

\subsection{Idea central}

Cuando una señal continua $f(t)$ se muestrea cada $T_s$ segundos
(es decir, con una frecuencia de muestreo $f_s = 1/T_s$), solo se
conservan los valores de $f(t)$ en esos instantes. Si la señal
original contiene componentes de frecuencia demasiado altas en
comparación con $f_s$, esos componentes de alta frecuencia
\textbf{no se pueden distinguir} de componentes de una frecuencia
mucho más baja: los datos muestreados de ambas señales resultan
ser idénticos. A esa frecuencia "falsa" de baja frecuencia que
aparece en su lugar se le llama \emph{alias}.

\subsection{Teorema de muestreo de Nyquist-Shannon}

Para evitar el aliasing, la frecuencia de muestreo debe ser mayor
al doble de la frecuencia máxima presente en la señal:

\begin{equation}
\boxed{\; f_s > 2 f_{\max} \;}
\end{equation}

A la frecuencia $2f_{\max}$ se le conoce como \textbf{frecuencia de
Nyquist}. Si esta condición no se cumple, las componentes de
frecuencia $f > f_s/2$ se "reflejan" hacia el rango
$[0, f_s/2]$, apareciendo como frecuencias falsas indistinguibles
de señales reales de menor frecuencia.

\subsection{Ejemplo clásico}

Un ejemplo cotidiano es el de las llantas de un automóvil o las
aspas de un helicóptero en una película: a veces parecen girar
muy lento, detenidas, o incluso en reversa, aunque en realidad
giran rápido. Esto sucede porque la cámara "muestrea" la escena
tomando una cantidad limitada de fotogramas por segundo; si esa
tasa es menor al doble de la velocidad de giro, el ojo humano
percibe un movimiento distinto (más lento o invertido) al
movimiento real, que es justamente un caso de aliasing.

\subsection{Relación práctica}

En electrónica y procesamiento digital de señales, este fenómeno
se evita colocando un \textbf{filtro anti-aliasing} (un filtro
pasa-bajas) antes del proceso de muestreo, para eliminar toda
componente de frecuencia mayor a $f_s/2$ y así garantizar que la
señal digitalizada represente fielmente a la señal original.

\section{Diferencia clave entre ambos fenómenos}

\begin{itemize}
    \item El \textbf{fenómeno de Gibbs} ocurre al truncar una serie
    de Fourier de una función \emph{continua en el dominio del
    tiempo}: el problema es cuántos armónicos se suman.
    \item El \textbf{fenómeno de aliasing} ocurre al
    \emph{muestrear} una señal en el tiempo: el problema es qué
    tan seguido se toman las muestras.
\end{itemize}

Ambos son limitaciones fundamentales que aparecen constantemente
en el análisis de Fourier y en el procesamiento digital de
señales, y es útil reconocerlos por separado para no confundirlos.

\end{document}
